% theme: animal_classification
\MajorQuestion{動物のなかま}
\SetRowSpaceQanda

\Instruction{背骨をもつ動物の分類に関する次の問いに答えなさい。}
\QQ{背骨をもつ動物を\NamedBlank{vertebrate}という。}{脊椎動物}
\QQ{\RefBlank{vertebrate}のうち、一生を水中で生活し、体表がうろこでおおわれるものを\NamedBlank{fish}という。}{魚類}
\QQ{\RefBlank{vertebrate}のうち、幼生は水中、成体は陸上や水辺で生活するものを\NamedBlank{amphibian}という。}{両生類}
\QQ{\RefBlank{vertebrate}のうち、体表がうろこでおおわれ、陸上に殻のある卵を産むものを何というか。}{は虫類}
\QQ{\RefBlank{vertebrate}のうち、体表が羽毛でおおわれるものを\NamedBlank{bird}という。}{鳥類}

\Instruction{呼吸と子の生まれ方に関する次の問いに答えなさい。}
\QQ{\RefBlank{vertebrate}のうち、体表が毛でおおわれ、子を母体内で育てるものを\NamedBlank{mammal}という。}{哺乳類}
\QQ{\RefBlank{fish}が呼吸に使う器官を何というか。}{えら}
\QQ{陸上で生活する\RefBlank{vertebrate}すべてに共通する呼吸器官は何か。}{肺}
\QQ{雌が卵を産み、卵から子が生まれる生まれ方を何というか。}{卵生}
\QQ{子が雌の体内である程度育ってから生まれる生まれ方を何というか。}{胎生}

\Instruction{体温と背骨をもたない動物に関する次の問いに答えなさい。}
\QQ{\RefBlank{bird}や\RefBlank{mammal}のように、体温をほぼ一定に保つ動物を何というか。}{恒温動物}
\QQ{外界の温度変化にともなって体温が変わる動物を何というか。}{変温動物}
\QQ{背骨をもたない動物を\NamedBlank{invertebrate}という。}{無脊椎動物}
\QQ{\RefBlank{invertebrate}のうち、体やあしに節があり、かたい殻をもつものを\NamedBlank{arthropod}という。}{節足動物}
\QQ{\RefBlank{arthropod}の体をおおう、かたい殻を何というか。}{外骨格}

\Instruction{節足動物と軟体動物に関する次の問いに答えなさい。}
\QQ{頭部・胸部・腹部に分かれ、胸部に3対のあしをもつ\RefBlank{arthropod}を\NamedBlank{insect}という。}{昆虫類}
\QQ{カニやエビなどが属する\RefBlank{arthropod}のなかまを何というか。}{甲殻類}
\QQ{イカやアサリなど、体とあしに節がなく、内臓が膜で包まれる動物を\NamedBlank{mollusk}という。}{軟体動物}
\QQ{\RefBlank{mollusk}の内臓を包む、筋肉でできた膜を何というか。}{外とう膜}
\QQ{\RefBlank{insect}が空気を取り入れる、胸部や腹部にある穴を何というか。}{気門}
